\section{Signed arithmetic operations in the exact original metrics}
The original polynomial and its eight marked inverse states are those of
FS13 and OE17, with the source attribution \cite{ES488,Tao}. The arithmetic
coordinates, including the four labels and every derivative unit, are
ISR1--12. OE1--8 applies the Type I representation of
Elsholtz and Tao \cite[Section~2]{ElsholtzTao2015}; it proves that the
original coefficient $A=-1/S$ is a unit for every positive witness in
the stated prime class. Here all signed operations are returned to the
literal original receivers, with their original moments and Gamma masses
\cite{WCF}. This calculation also records how their action on functions
produces the original nonlinear signed states.

Write $v(r)=(1,r,r^2,r^3)^{\mathsf T}$, $V_{ij}=r_i^j$,
$f_i=f/(r-r_i)$, $d_i=f'(r_i)$, and let $c_i$ be the full coefficient
column of $f_i$. Thus $v(r_i)^{\mathsf T}c_i=d_i$ and $d_i\ne0$.
For $\epsilon\in\{1,-1\}^4$ put
\[
 B_\epsilon=V^{-1}\operatorname{diag}(\epsilon_i)V,
 \qquad \Sigma_\epsilon=\operatorname{diag}(I_4,B_\epsilon).
 \tag{SGR1}
\]
In the original parity basis, this map fixes $R(r)$ and sends
$\sigma Q(r)$ to $\sigma\eta_\epsilon(r)Q(r)$, where
$\eta_\epsilon=\sum_i\epsilon_i f_i/d_i$ modulo $f$.
Evaluation gives $\eta_\epsilon(r_i)=\epsilon_i$; hence
$\eta_\epsilon^2=1$ modulo $f$ and $\Sigma_\epsilon^2=I_8$.
This proves the algebra automorphism over the generic field, including
its inverse and its exact sign convention.

\subsection{The complete integral intersection and invariant suborder}
Set $R=\mathbb Z_p$, $E=R[r]\oplus\sigma R[r]$ modulo the original
relations. For a single sign change at $r_j$, the odd block is
\[
 B_j=I_4-2c_jv(r_j)^{\mathsf T}/d_j,
 \qquad
 E\cap\Sigma_jE
 =\{R_0+\sigma Q:Q(r_j)\in d_jR\}.
 \tag{SGR2}
\]
Indeed $c_j$ is a primitive column because its last entry is one.
Since two is a unit, the expression
$Q-2c_jQ(r_j)/d_j$ is integral if and only if $Q(r_j)/d_j$ is
integral. Division by the monic $r-r_j$ proves the following complete
eight-column basis of that intersection:
\[
 D_j=\operatorname{diag}(I_4,N_j),\qquad
 N_j=\begin{pmatrix}
 d_j&-r_j&0&0\\0&1&-r_j&0\\0&0&1&-r_j\\0&0&0&1
 \end{pmatrix},\qquad
 E/(E\cap\Sigma_jE)\simeq R/(d_j).
 \tag{SGR3}
\]
The quotient map is exactly $R_0+\sigma Q\mapsto Q(r_j)\bmod d_j$;
it is onto by testing the constant polynomial. Both the unit and the
valuation of $d_j$ are retained in $N_j$.

If $b_j=v_p(d_j)>0$, the relative Smith exponents of $B_j$ are
$(-b_j,0,0,b_j)$. To prove this, choose an integral invertible basis
matrix whose first column is $c_j$. Such a matrix exists because
$c_j$ is primitive. The row $v(r_j)^{\mathsf T}$ in this basis has
first entry $d_j$ and remains primitive, so one of its remaining
entries is a unit. Integral operations on these other three columns
give a row $(d_j,u,0,0)$ with $u$ a unit. The resulting nontrivial
two-by-two block is
$\left(\begin{smallmatrix}-1&-2u/d_j\\0&1\end{smallmatrix}\right)$.
Its least entry valuation is $-b_j$ and its determinant is $-1$;
integral row and column elimination therefore gives exponents
$-b_j,b_j$. The other two entries are units. If $b_j=0$, the whole
matrix and its inverse are integral and all exponents are zero.
The four even exponents in $\Sigma_j$ are always zero.

The same formula covers every nonconstant sign assignment on an actual
residue cluster: multiply by the sign assignment constant on that
cluster to leave a single exceptional sign. Cluster-constant signs
act integrally. To prove that assertion without a residue-only
argument, factor $f$ into its residue clusters. The resultants between
distinct clusters are units, so the Bezout inverse, obtained by solving
the Sylvester system with its unit determinant, lies in $R[r]$.
The resulting exact Chinese-remainder idempotents give the required
integral sign polynomial and its integral inverse. Conversely, an
integral polynomial has equal residues at two roots congruent modulo
$p$; signs differing by two cannot be its values at those roots.
Thus the integral sign subgroup consists exactly of the signs constant
on each cluster: its order is eight in Type I and four in Type II.
For a mixed pair or triple cluster SGR3 gives respective cyclic
defects $R/(p)$ and $R/(p^2)$, with all full units still in the map.

Let $Q_f$ denote the complete matrix of columns $c_i$, as in ISR12.
The largest suborder of $E$ stable under all sixteen signed operations is
\[
 J=\bigcap_\epsilon\Sigma_\epsilon E
   =R[r]\oplus\sigma Q_fR^4,
 \qquad
 E/J\simeq R^4/Q_fR^4.
 \tag{SGR4}
\]
Testing every single sign change proves that membership is equivalent
to $Q(r_i)/d_i\in R$ for every $i$. The complete Lagrange formula
$Q=\sum_i Q(r_i)f_i/d_i$ proves the displayed equality. The
intersection is a subring because every $\Sigma_\epsilon E$ is a
subring; it is stable because the sign group permutes its factors.
Every stable suborder inside $E$ is contained in this intersection,
proving maximality. Finally $Q_f=H_f^{-1}V^{\mathsf T}$, using
ISR9--10 and $VQ_f=\operatorname{diag}(d_i)$. The integral matrix
$H_f^{-1}$ is invertible, so the Smith exponents of $E/J$ are those
of $V$: $(0,0,0,1)$ in Type I and $(0,0,1,2)$ in Type II.
One obtains these lists by using the exact integral cluster
decomposition above and, in each size-$n$ cluster, the successive
Newton columns $1,(r-r_0),\ldots,\prod_{a<n-1}(r-r_a)$.
All within-cluster gaps have valuation one; after dividing the $j$th
Newton column by $p^j$, its evaluation matrix has unit diagonal and
is invertible over $R$. This proves the asserted factors without
discarding the full gap units.

\subsection{Original affine states as exact functions of the signed algebra}
In the generic algebra let $D(r)=Af'(r)\bmod f$. Its inverse exists
and is exactly $D^{-1}=\sum_i D_i^{-1}f_i/d_i$. In each of the
following expressions polynomial remainders are taken modulo the
unchanged monic $f$. Define the four original coordinate functions by
the full parity columns
\[
 C_{\rm aff}=\begin{pmatrix}
 0&-r&3iD&(B-17r+Ar^2)D-2AD^2\\
 D^{-1}&-i&AD+2r&7ir^2-13iD
 \end{pmatrix}_{\!\text{coefficient columns}}.
 \tag{SGR5}
\]
Here each displayed entry is its four-entry coefficient column, so
$C_{\rm aff}$ is an $8$-by-$4$ matrix. For either $s^2=D_i$ set
$e_{i,s}=(v(r_i)^{\mathsf T},s v(r_i)^{\mathsf T})$. Direct substitution
proves, retaining every term in FS13,
\[
 q_{i,s}=C_{\rm aff}^{\mathsf T}e_{i,s}^{\mathsf T},\qquad
 e_{i,s}\Sigma_\epsilon=e_{i,\epsilon_i s},\qquad
 C_{\rm aff}^{\mathsf T}\Sigma_\epsilon^{\mathsf T}e_{i,s}^{\mathsf T}
   =q_{i,\epsilon_i s}.
 \tag{SGR6}
\]
For example $s/D_i=s^{-1}$, $sD_i=s^3$, and $D_i^2=s^4$,
which give the first, third and fourth original affine coordinates
with their precise signs and constants. The middle equality is a
row-by-row matrix product using $v(r_i)^{\mathsf T}B_\epsilon
=\epsilon_i v(r_i)^{\mathsf T}$. It states the dual action explicitly.
Thus the function-algebra action and the signed affine states are
related by the exact displayed morphism; they are not identified
merely because they have the same labels.

For the original four-dimensional receiving maps $\Phi,\Psi$ with
$T_{\mathcal A}\Phi=\Psi$, the corresponding original vectors are
$\Phi q_{i,s}$ and $\Psi q_{i,s}$. Equation SGR6 followed by
$\Phi$ or $\Psi$ gives their complete signed-state return, and
\[
 T_{\mathcal A}\Phi C_{\rm aff}^{\mathsf T}
       \Sigma_\epsilon^{\mathsf T}e_{i,s}^{\mathsf T}
 =\Psi C_{\rm aff}^{\mathsf T}e_{i,\epsilon_i s}^{\mathsf T}.
 \tag{SGR7}
\]
This is a pointwise exact relation on all eight original states.
The complete rational sign involution on its affine domain remains
the nonlinear map already derived in the manuscript; SGR7 does not
replace it by a linear operator on the four affine coordinates.

\subsection{Every original-metric signed-operation singular value}
Write $\mathcal F_b=F_b\oplus F_b$ and
$\Gamma_b=\operatorname{diag}(H_b,H_b)$ for the unchanged receivers
and their coefficient Grams in GD14 and ECC24--28, $b=U,W$.
Define the actual transported operation
\[
 \mathscr S_{b,\epsilon}=\mathcal F_b\Sigma_\epsilon\mathcal F_b^{-1},
 \qquad
 (T_{\mathcal A}\oplus T_{\mathcal A})\mathscr S_{U,\epsilon}
   =\mathscr S_{W,\epsilon}(T_{\mathcal A}\oplus T_{\mathcal A}).
 \tag{SGR8}
\]
The equality follows by substitution of
$(T_{\mathcal A}\oplus T_{\mathcal A})\mathcal F_U=\mathcal F_W$;
all maps have the displayed domains and inverses. The squared
singular values are exactly the eigenvalues of
$\Gamma_b^{-1}\Sigma_\epsilon^*\Gamma_b\Sigma_\epsilon$.
Here is their full finite evaluation in terms of the original $H_b$.

For a single sign change put
\[
 \nu_{b,j}=
 \frac{(c_j^*H_bc_j)
        (v(r_j)^{\mathsf T}H_b^{-1}\overline{v(r_j)})}{|d_j|^2},
 \qquad s_{b,j}=\sqrt{\nu_{b,j}}+\sqrt{\nu_{b,j}-1}.
 \tag{SGR9}
\]
The dual Cauchy--Schwarz inequality and $v(r_j)^{\mathsf T}c_j=d_j$
give $\nu_{b,j}\ge1$. In isometric coordinates the rank-one
projection $c_jv(r_j)^{\mathsf T}/d_j$ becomes $uw^*$ with
$w^*u=1$ and $\|u\|^2\|w\|^2=\nu_{b,j}$. On the span of $u,w$
the reflection $I-2uw^*$ has determinant of absolute value one.
Its squared singular values have product one and sum
$4\nu_{b,j}-2$, obtained by expanding its two-by-two trace.
Solving that quadratic gives $s_{b,j}^2,s_{b,j}^{-2}$.
The orthogonal complement is fixed. This also covers $\nu=1$,
when the projection is orthogonal. Hence the full eight singular
values and all nontrivial exterior norms are
\[
 (s_{b,j},1,1,1,1,1,1,s_{b,j}^{-1}),\qquad
 \|\wedge^k\mathscr S_{b,j}\|=
 \begin{cases}s_{b,j}&1\le k\le7,\\1&k=8.\end{cases}
 \tag{SGR10}
\]

For any sign pattern, change only the coefficient description to its
full evaluation basis and put $K_b=(V^{-1})^*H_bV^{-1}$.
Order the positive and negative labels, retaining the permutation, and
write $K_b=\left(\begin{smallmatrix}A_b&C_b\\C_b^*&B_b\end{smallmatrix}\right)$.
If either label set is empty, the signed operation is an isometry.
Otherwise both diagonal blocks are positive. Let $\rho_{b,a}$ be
all singular values of $A_b^{-1/2}C_bB_b^{-1/2}$, including zeros,
for $1\le a\le m=\min(n_+,n_-)$, in decreasing order. Then
\[
 0\le\rho_{b,a}<1,\qquad
 s_{b,a}=\sqrt{\frac{1+\rho_{b,a}}{1-\rho_{b,a}}},\qquad
 \operatorname{Sing}(\mathscr S_{b,\epsilon})
   =(s_{b,1},\ldots,s_{b,m},
          \underbrace{1,\ldots,1}_{8-2m},
          s_{b,m}^{-1},\ldots,s_{b,1}^{-1}).
 \tag{SGR11}
\]
These formulas retain every entry of the original metric. Positive
definiteness of $K_b$ proves $I-W^*W>0$ for the displayed cross
matrix $W$, and therefore $\rho<1$. Its singular-value decomposition
reduces the metric to two-by-two blocks
$\left(\begin{smallmatrix}1&\rho\\\rho&1\end{smallmatrix}\right)$
and unmatched unit blocks. The sign reflection changes the two
off-diagonal entries to $-\rho$. The generalized squared eigenvalues
are $(1+\rho)/(1-\rho)$ and its reciprocal, proving SGR11.
In particular the exact $k$th exterior norm is
\[
 \|\wedge^k\mathscr S_{b,\epsilon}\|
    =\prod_{a=1}^{\min(k,m,8-k)}s_{b,a}\quad(0\le k\le8),
 \tag{SGR12}
\]
with an empty product equal to one. This includes every sign pattern,
including the two-sign pattern of a nontrivial unramified quadratic
action, without assuming that its integral baseline is unitary.

\subsection{The intersection lattice in both original receivers}
The passage between arithmetic and complex norms uses a common rational
model, not a tacit embedding of $\mathbb Q_p$ into $\mathbb C$.
For the actual integer roots set $R_{(p)}=\mathbb Z_{(p)}$.
The coefficients $A,f,D_j,Q_f,\Sigma_\epsilon$ lie in $\mathbb Q$,
and the defining order and its intersection are defined over $R_{(p)}$.
The completion $R_{(p)}\to\mathbb Z_p$ gives exactly SGR2--4,
while $R_{(p)}\to\mathbb Q\to\mathbb C$ gives their complex
coefficient spaces. The full comparison is the base change of
\[
 0\longrightarrow D_jR_{(p)}^8\longrightarrow R_{(p)}^8
 \xrightarrow{\;R_0+\sigma Q\mapsto Q(r_j)\bmod d_j\;}
 R_{(p)}/(d_j)\longrightarrow0.
 \tag{SGR15}
\]
Exactness is the same monic-division proof as SGR3. Smith elimination
over $R_{(p)}$, a discrete valuation ring, proves that completion
retains the cyclic factor $p^{b_j}$. Complex base change is exact:
localization to $\mathbb Q$ is flat, and extension from the field
$\mathbb Q$ to $\mathbb C$ is flat. It kills the finite cyclic
factor because $d_j\ne0$ is invertible in $\mathbb C$.
The functions in SGR5 are then defined over the same rational model
after adjoining the displayed $i$, and are evaluated over $\mathbb C$.
This is the exact arithmetic-to-complex correspondence used here;
it selects no map of topological fields from $\mathbb Q_p$ to
$\mathbb C$. The same common-model construction applies to the
rational generator matrices in OE18 and PS23--25.

For the full basis $D_j$ in SGR3, define
$\Gamma_{b,j}=D_j^*\Gamma_bD_j$ and the images
$\mathcal F_bD_jR_{(p)}^8$ in the corresponding complex receiver.
Both complete generator maps, including the same cyclic defect, are
\[
 (T_{\mathcal A}\oplus T_{\mathcal A})\mathcal F_UD_j
       =\mathcal F_WD_j,\qquad
 \det\Gamma_{b,j}=|d_j|^2\det\Gamma_b.
 \tag{SGR13}
\]
The complex spans of these lattices have the full original dimension
eight. Their identity-coefficient conductor has squared inverse
singular matrix $\Gamma_{W,j}^{-1}\Gamma_{U,j}$, which satisfies
the exact similarity
\[
 \Gamma_{W,j}^{-1}\Gamma_{U,j}
     =D_j^{-1}\Gamma_W^{-1}\Gamma_UD_j.
 \tag{SGR14}
\]
Consequently every original inverse singular value and every inverse
exterior norm is unchanged by this simultaneous lattice restriction
with its actual source and target metrics. This is a full-spectrum
identity, not a determinant-only cancellation. The same proof applies
to $J$ using $\operatorname{diag}(I_4,Q_f)$, and to a mixed sign
assignment using its exact common intersection basis.
Equations SGR8--14 give the complete operator, metric and exterior
return of the arithmetic signed operation. The finite arithmetic
failure to preserve $E$ is measured by SGR3 and does not by itself
produce a singularity of the original complex linear conductor.
None of these calculations identifies a local quadratic action with
the arithmetic operator required for an RH conclusion.
